\chapter*{Abstract}
\addcontentsline{toc}{chapter}{Abstract}

Rehearsal-free class-incremental learning with parameter-efficient
fine-tuning asks how a low-rank-adapted model can absorb a sequence of new
classes without storing past data and without catastrophically forgetting
what it already knows. This thesis studies two ways of organising LoRA
adaptation across such a sequence on a frozen CLIP ViT-B/16 backbone.
\emph{SimpleAvg} trains an independent, fixed-rank (rank-80) LoRA
specialist for each step and merges all specialists once, at the end, into
a single dense weight update; it is treated throughout as a strong
reference baseline rather than as a proposed contribution. \emph{RankExt}
instead grows one persistent, cumulative low-rank adapter across the
incremental sequence, adding a new rank block at each step while all earlier blocks
remain frozen but active, so that adaptation capacity is introduced
incrementally rather than specialist by specialist. Both strategies are
evaluated, with and without knowledge distillation and orthogonality-based
stabilisation, under a five-step/twenty-class-per-step class-incremental
protocol on CIFAR-100 as the primary benchmark, with CUB-200-2011
($5\times40$) as a second-dataset supporting evaluation, using paired
task-agnostic (open)
and task-oracle (restricted) evaluation to help distinguish what a model
still discriminates from what it can still express under output-space
competition with later steps.

SimpleAvg is unexpectedly strong on its own, reaching 75.16\% all-seen
accuracy with no auxiliary mechanism. Plain RankExt is far weaker, at
34.99\%, despite retaining high restricted accuracy, suggesting that
cross-step output-space competition contributes substantially to its poor
open performance while within-step discrimination remains comparatively
strong. This divergence between
restricted and open performance, which persists even where open accuracy
is poor, is a central diagnostic finding of this thesis. Auxiliary
stabilisation recovers most of RankExt's gap: combining knowledge
distillation with projected feature protection and a factor-space
orthogonality penalty brings its highest observed configuration to
70.76\% all-seen accuracy, 4.45 percentage points below the highest observed
SimpleAvg result (75.21\%). On CUB-200-2011 the same qualitative pattern
appears: plain SimpleAvg reaches 69.47\%, and the highest observed RankExt
configuration 63.43\%. The
comparison is not framed as a single winner-take-all contest: the central
question is how closely a persistent strategy that grows its adaptation
capacity incrementally can approach a strong, independently-trained
reference, and at what retention--plasticity cost. All headline benchmark
results are single-seed observations; multi-seed replication is identified
as a priority for confirming the findings.
